% Downstream language-modelling evaluation. Self-contained so it survives a rewrite of
% the surrounding paper: \input this file wherever the Results section lives.
%
% Requires from the preamble:
%   \newcommand{\bnd}[1]{\texttt{bnd\_#1}}
%   \usepackage{booktabs}
% Requires in the bibliography:
%   li2024datacomplm   (DCLM, for the CORE metric)
% Paths assume the main .tex sits in marker_experiments/paper/.
% (LaTeX resolves \input relative to the main document, not to this file).
% Reads: generated/downstream_tables.tex, which make_tex_tables.py writes from
%   marker_experiments/downstream/{manifest,results,text_stats}.* and must not be hand-edited.
%
\subsection{Downstream language modelling}
\label{sec:downstream}

Compression is not quality, so we pretrain a model on each vocabulary. Arms are matched on
\emph{total} vocabulary rather than on merges added: the grids above fix
\texttt{additional\_vocab\_size}, which leaves the marker arms a few atomic tokens larger, and
that difference should not be read as a modelling effect. Every arm is retrained at a total
vocabulary of 34{,}685 on 5\,GB of FineWeb English. Table~\ref{tab:downstream-compression} reports the
resulting tokenizers. The gain, $+2.87\%$, is smaller than the $+3.77\%$ of
Table~\ref{tab:main}. Three things differ between the two: the tokenizer-training corpus
(5\,GB of FineWeb here against 1\,GB of FineWiki there), the vocabulary convention, and the
fact that chars/token is measured on a held-out FineWiki slice in both cases, which is out of
domain here. We do not separate them. Holding total vocabulary fixed accounts for almost none
of it: it costs \bnd{wpd} one merge in 33{,}000 relative to \texttt{plain}.

The gap between corpora matters for reading the result. Measured on ClimbMix, the corpus the
models actually read, the same tokenizers gain $+0.59\%$ (\bnd{wpd}) and $+0.70\%$
(\bnd{wpd\_caps}) chars/token over \texttt{plain}, and \bnd{w} loses $7.78\%$ rather than
$14.85\%$. The compression advantage the models could exploit is therefore several times
smaller than the headline FineWiki figure.

Models are nanochat at depth 12 on ClimbMix, scored on validation bits-per-byte and on DCLM
CORE \citep{li2024datacomplm}, in Table~\ref{tab:downstream-lm}.

% Generated by marker_experiments/downstream/make_tex_tables.py. Do not edit.
% sources: marker_experiments/downstream/manifest.json, results/marker_downstream/results.tsv

\begin{table}[t]
\centering
\small
\begin{tabular}{lrrrr}
\toprule
Arm & Vocab & chars/token & $\Delta$ (\%) & RT fail \\
\midrule
plain & 34{,}685 & 3.6360 & --- & 0 \\
\bnd{w} & 34{,}685 & 3.0960 & $-14.85$ & 0 \\
\bnd{wpd} & 34{,}685 & 3.7403 & $+2.87$ & 0 \\
\bnd{wpd\_caps} & 34{,}685 & 3.7456 & $+3.02$ & 0 \\
\bottomrule
\end{tabular}
\caption{Vocabulary-matched tokenizers on fineweb\_en\_5gb, all at 34{,}685 total vocabulary. Chars/token on the held-out FineWiki English slice (500 documents, 3{,}602{,}925 characters); $\Delta$ relative to \texttt{plain}. RT fail is roundtrip failures.}
\label{tab:downstream-compression}
\end{table}

\begin{table}[t]
\centering
\small
\begin{tabular}{lrrr}
\toprule
Arm & $n$ & val bpb & CORE \\
\midrule
\multicolumn{4}{c}{\textit{no downstream runs finished yet}} \\
\bottomrule
\end{tabular}
\caption{nanochat depth-12 pretraining on ClimbMix. $n$ is runs averaged; $\pm$ is the standard deviation across seeds. CORE is \textit{n/a} for the arms whose tokenization does not satisfy the prefix property its language-modelling tasks assume (\S\ref{sec:downstream}).}
\label{tab:downstream-lm}
\end{table}


All three boundary arms reach lower bits-per-byte than the baseline. Because the seeds are
paired across arms, the differences are paired: \texttt{plain} minus \bnd{w} is
$+0.00853 \pm 0.00058$, minus \bnd{wpd\_caps} $+0.00578 \pm 0.00043$, and minus \bnd{wpd}
$+0.00532 \pm 0.00055$, each 16 to 25 times its standard deviation over three seeds.
\bnd{wpd} and \bnd{wpd\_caps} are indistinguishable ($+0.00046 \pm 0.00087$), so the two
shift codes cost nothing measurable.

Two qualifications. \bnd{w} appears strongest but emits $8.4\%$ more tokens per byte than
\texttt{plain}, so its model spends proportionally more forward passes on the same text and
the metric does not correct for that; \bnd{wpd} and \bnd{wpd\_caps} emit $0.6$ to $0.7\%$
\emph{fewer}, so their improvement cannot come from that source and they are the evidence
that carries the claim. And the improvement is close in size to the compression gain on the
corpus the models read: $+0.59\%$ chars/token for \bnd{wpd} on ClimbMix against $+0.60\%$
bits-per-byte.

\paragraph{CORE cannot score the punctuation-marking arms.} Its language-modelling tasks
tokenize a context and that context plus its continuation, and require the first to be a
prefix of the second. A marker that depends on the following character does not satisfy this:
under \bnd{wpd}, \texttt{':'} is one token at the end of a string and a different token when a
space follows, so extending \texttt{"Answer:"} with \texttt{" Boston"} rewrites a token
already emitted. Across the nine CORE language-modelling tasks, 512 examples at up to 60 per
task, this affects 0/512 examples for \texttt{plain} and \bnd{w} and 399/512 for \bnd{wpd}
and \bnd{wpd\_caps}. The counts are identical for the Korean- and Russian-trained tokenizers
on the same English prompts, so they follow from the marker scheme rather than from the
tokenizer's training language. On the compression-grid tokenizers, which are trained at a
different vocabulary and corpus and are not among the arms evaluated here, \bnd{wp} gives
311/512, placing it between the two extremes. We report CORE
only where it is defined, and bits-per-byte everywhere. This is a property of the evaluation
protocol rather than of the models, but it constrains what any context-dependent marker scheme
can be compared on.
